% !TeX root = ../main.tex

\begin{nomenclature}

\begin{longtable}{r l}
  \multicolumn{2}{l}{\textbf{符号与缩写}} \\[0.5em]

  % 基本物理量
  $T$ & 温度 [K] \\
  $P$ & 功率 [W] \\
  $I$ & 电流 [A] \\
  $U$ & 电压 [V] \\
  $n$ & 物质的量 [mol] \\
  $V$ & 体积 [$\mathrm{m}^3$] \\
  $v$ & 体积流量 [$\mathrm{m}^3/\mathrm{s}$] \\

  \multicolumn{2}{l}{} \\[-0.3em]
  \multicolumn{2}{l}{\textbf{状态变量}} \\[0.3em]
  $T_{s,in}$ & 换热器碱液出口温度 [K] \\
  $T_{s,i}$ & 第$i$个电解槽温度 [K], $i=1,2,3,4$ \\
  $T_{sep}$ & 分离器温度 [K] \\
  $T_{c,out}$ & 冷却水出口温度 [K] \\
  $n_{H_2,an,i}$ & 第$i$个阳极室溶解氢气摩尔数 [mol] \\
  $n_{liq}$ & 分离器液相氢气摩尔数 [mol] \\
  $n_{gas}$ & 分离器气相氢气摩尔数 [mol] \\

  \multicolumn{2}{l}{} \\[-0.3em]
  \multicolumn{2}{l}{\textbf{控制输入}} \\[0.3em]
  $I_i$ & 第$i$个电解槽电流 [A] \\
  $v_{lye,i}$ & 第$i$个电解槽碱液流量 [$\mathrm{m}^3/\mathrm{s}$] \\
  $v_c$ & 冷却水流量 [$\mathrm{m}^3/\mathrm{s}$] \\

  \multicolumn{2}{l}{} \\[-0.3em]
  \multicolumn{2}{l}{\textbf{系统参数}} \\[0.3em]
  $N$ & 电解槽数量 (本研究$N=4$) \\
  $N_{cell}$ & 单槽电芯数 (本研究$N_{cell}=368$) \\
  $A_{cell}$ & 电芯面积 [$\mathrm{m}^2$] \\
  $P_{sys}$ & 系统压力 [Pa] \\
  $F$ & 法拉第常数 [C/mol] \\
  $R$ & 理想气体常数 [J/(mol$\cdot$K)] \\

  \multicolumn{2}{l}{} \\[-0.3em]
  \multicolumn{2}{l}{\textbf{评价指标}} \\[0.3em]
  HTO & 氢气至氧侧传输量百分比 (Hydrogen Transport to Oxygen) [\%] \\
  RMSE & 均方根误差 (Root Mean Square Error) \\
  MAE & 平均绝对误差 (Mean Absolute Error) \\

  \multicolumn{2}{l}{} \\[-0.3em]
  \multicolumn{2}{l}{\textbf{方法与算法}} \\[0.3em]
  AWE & 碱性电解水 (Alkaline Water Electrolysis) \\
  NMPC & 非线性模型预测控制 (Nonlinear MPC) \\
  DDPM & 去噪扩散概率模型 (Denoising Diffusion Probabilistic Model) \\
  TCN & 时序卷积网络 (Temporal Convolutional Network) \\
  MLP & 多层感知机 (Multi-Layer Perceptron) \\
  CBF & 控制障碍函数 (Control Barrier Function) \\

  \multicolumn{2}{l}{} \\[-0.3em]
  \multicolumn{2}{l}{\textbf{下标与上标}} \\[0.3em]
  $ref$ & 参考值 \\
  $min/max$ & 最小/最大值 \\
  $in/out$ & 入口/出口 \\
  $lye$ & 碱液 \\
  $cw$ & 冷却水 \\
  $amb$ & 环境 \\

\end{longtable}

\end{nomenclature}
